\documentclass[aspectratio=169]{beamer}
\usetheme{SimplePlus}

\usepackage{tikz}
\usepackage{amsmath}
\usepackage{amssymb}
\usepackage{bm}

\title[Traces and Poincar\'e]{Applications of Functional Analysis and PDEs: \\ Traces and Poincar\'e Inequalities}
\author{Nicol\'as A. Barnafi}
\date{}

\begin{document}

\begin{frame}
    \titlepage
\end{frame}

\begin{frame}{Outline}
    \tableofcontents
\end{frame}

\section{Traces}

\begin{frame}{The Boundary Value Problem}
    To rigorously define boundary conditions for PDEs, we must restrict functions in $\mathbb{R}^n$ to the boundary $\partial\Omega \subset \mathbb{R}^{n-1}$[cite: 662].

    \vspace{0.2cm}
    \textbf{The Challenge:}
    The boundary $\partial\Omega$ has Lebesgue measure $0$ in $\mathbb{R}^n$[cite: 664]. Functions in Sobolev spaces are equivalence classes defined \textit{almost everywhere}. Pointwise evaluation on a set of measure zero is meaningless.

    \vspace{0.2cm}
    \textbf{The Solution:}
    We need trace operators that require a certain level of regularity to guarantee the restriction makes sense[cite: 664].
\end{frame}

\begin{frame}{The Dirichlet Trace}
    Let $\Omega \subset \mathbb{R}^n$ be a bounded, Lipschitz domain[cite: 676].

    We start by defining the \textbf{trace operator} $\gamma_0$ for smooth functions[cite: 676]:
    \[
        \gamma_0 : \mathcal{D}(\bar{\Omega}) \to C^\infty(\partial\Omega), \quad \gamma_0 u := u|_{\partial\Omega} \text{}
    \]

    \textbf{Trace Inequality:} There exists $C > 0$ such that[cite: 679, 680]:
    \[
        \|\gamma_0 f\|_{0,\partial\Omega} \le C\|f\|_{1,\Omega} \quad \forall f \in \mathcal{D}(\bar{\Omega}) \text{}
    \]

    \textbf{Trace Theorem:} By density, $\gamma_0$ extends to a continuous and surjective operator[cite: 685, 700]:
    \[
        \gamma_0 : H^1(\Omega) \to H^{1/2}(\partial\Omega) \text{}
    \]
\end{frame}

\begin{frame}{Spaces with Boundary Conditions}
    We can now formally define Sobolev spaces with zero boundary conditions via the kernel of the trace operator[cite: 696, 697]:
    \[
        H^1_0(\Omega) := \ker \gamma_0 = \{ u \in H^1(\Omega) : \gamma_0 u = 0 \} \text{}
    \]
    \begin{itemize}
        \item The dual space of $H^1_0(\Omega)$ is denoted $H^{-1}(\Omega)$[cite: 699].
        \item The dual space of the trace space $H^{1/2}(\partial\Omega)$ is $H^{-1/2}(\partial\Omega)$[cite: 702].
    \end{itemize}

    \vspace{0.3cm}
    \textit{Formal Interpretation:} The Dirichlet condition $u = g$ on $\partial\Omega$ for the Poisson problem makes sense if and only if $g \in H^{1/2}(\partial\Omega)$, due to the surjectivity of the trace operator[cite: 706].
\end{frame}

\begin{frame}{Normal and Tangential Traces}
    To handle Neumann conditions, we extend traces to $H(\text{div})$ and $H(\text{curl})$ using integration by parts (Divergence Theorem)[cite: 714, 715].

    \vspace{0.2cm}
    \textbf{Normal Trace Theorem ($H(\text{div})$)}[cite: 731, 732, 733]:
    The mapping $\gamma_N v = v|_{\partial\Omega} \cdot n$ extends to a continuous linear map $\gamma_N : H(\text{div}; \Omega) \to H^{-1/2}(\partial\Omega)$.
    \[
        \langle \gamma_N v, \phi \rangle_{-1/2, 1/2} = (v, \nabla \phi) + (\text{div } v, \phi) \quad \forall v \in H(\text{div}; \Omega), \phi \in H^1(\Omega) \text{}
    \]

    \vspace{0.2cm}
    \textbf{Tangential Trace Theorem ($H(\text{curl})$)}[cite: 752, 753]:
    The extension $\gamma_t v = n \times v$ maps continuously: $\gamma_t : H(\text{curl}; \Omega) \to H^{-1/2}(\partial\Omega)$.
    \[
        (\text{curl } v, \phi) - (v, \text{curl } \phi) = \langle \gamma_t v, \phi \rangle \quad \forall v \in H(\text{curl}; \Omega), \phi \in H^1(\Omega) \text{}
    \]
\end{frame}

\section{Poincar\'e Inequalities}

\begin{frame}{The Poincar\'e Inequality}
    \textbf{Lemma (Poincar\'e inequality):} Let $\Omega \subset \mathbb{R}^n$ be a bounded Lipschitz domain. There exists $C > 0$ such that[cite: 762, 763, 764]:
    \[
        \|u\|_{0,\Omega} \le C\|\nabla u\|_{0,\Omega} \quad \forall u \in H^1_0(\Omega) \text{}
    \]

    \vspace{0.3cm}
    \textit{Consequence:} The seminorm $|x|_{H^1} := \|\nabla x\|_{0,\Omega}$ is equivalent to the full $H^1$ norm for functions with homogeneous Dirichlet conditions[cite: 780, 781].

    \vspace{0.3cm}
    We will prove this using a powerful contradiction argument relying on \textbf{Compact Embeddings}.
\end{frame}

\begin{frame}{Proof via Compactness (Step 1: Setup Contradiction)}
    Assume the inequality is false.

    There exists a sequence $(v_n)_n \subset C^\infty_0(\Omega) \subset H^1_0(\Omega)$ such that[cite: 765, 766]:
    \[
        \|v_n\|_{0,\Omega} = 1 \quad \text{and} \quad \|\nabla v_n\|_{0,\Omega} < \frac{1}{n} \quad \forall n \in \mathbb{N} \text{}
    \]

    \textbf{Analysis of the sequence:}
    \begin{itemize}
        \item The full $H^1$ norm is $\|v_n\|_{1,\Omega} = \|v_n\|_{0,\Omega} + \|\nabla v_n\|_{0,\Omega} < 1 + \frac{1}{n}$.
        \item Therefore, $(v_n)_n$ is bounded in $H^1(\Omega)$[cite: 766].
    \end{itemize}
\end{frame}

\begin{frame}{Proof via Compactness (Step 2: Apply Embeddings)}
    Since $H^1(\Omega)$ is a reflexive space, the bounded sequence $(v_n)_n$ must have a \textbf{weakly convergent subsequence} $(v_{n_k})_k \rightharpoonup v \in H^1(\Omega)$[cite: 766].

    \vspace{0.3cm}
    \textbf{The Compactness Argument:}
    By the \textit{Rellich-Kondrachov Theorem}, the embedding $H^1(\Omega) \hookrightarrow \hookrightarrow L^2(\Omega)$ is compact[cite: 654, 766].
    \begin{itemize}
        \item Weak convergence in $H^1$ implies \textbf{strong convergence} in $L^2(\Omega)$[cite: 766].
        \item Therefore, $v_{n_k} \to v$ strongly in $L^2(\Omega)$.
        \item This preserves the norm exactly: $\|v\|_{0,\Omega} = \lim_{k\to\infty} \|v_{n_k}\|_{0,\Omega} = 1$[cite: 766].
    \end{itemize}
\end{frame}

\begin{frame}{Proof via Compactness (Step 3: The Contradiction)}
    Now consider the gradients. For any test vector field $\phi \in [C^\infty_0(\Omega)]^d$, weak convergence in $H^1$ gives[cite: 769, 770]:
    \[
        (\nabla v, \phi) = -(v, \nabla \cdot \phi) = -\lim_{k\to\infty} (v_{n_k}, \nabla \cdot \phi) = \lim_{k\to\infty} (\nabla v_{n_k}, \phi) \text{}
    \]
    But we assumed $\|\nabla v_{n_k}\|_0 < \frac{1}{n_k} \to 0$. Therefore, $(\nabla v, \phi) = 0$, meaning $\nabla v = 0$[cite: 770].

    \vspace{0.3cm}
    \begin{itemize}
        \item Since $\Omega$ is connected, $\nabla v = 0 \implies v = c$ (a constant) almost everywhere[cite: 770].
        \item Since $v_{n_k} \in H^1_0(\Omega)$, the trace evaluates to 0. Thus, $v|_{\partial\Omega} = c = 0$[cite: 771].
        \item This implies $v \equiv 0$, which contradicts our earlier finding that $\|v\|_{0,\Omega} = 1$[cite: 771].
    \end{itemize}
    \qed
\end{frame}

\begin{frame}{Poincaré Inequality}
    STATE THAT NORMS ARE EQUIV IN H01
    MENTION EQUIV IN NULL AVG
\end{frame}

\begin{frame}
    \maketitle
\end{frame}
\end{document}
